\scanpage{18}
\deleted{En plus, on \illegible{1} \unsure{considère} l'\unsure{homom.}}

Considérons \deleted{pour} $A\in C$, d'où \deleted{$\widetilde A\in\widetilde C$ et} $i_A:\deleted{\unsure{A\longrightarrow T\widetilde A}}\inserted{A\mathrel{\unsure{-|\to}}TSA}$ \deleted{\unsure{$\widetilde A$,}} et par suite $S(i_A):SA\longrightarrow STSA$ ; d'autre part, on a aussi $(ST)(SA)\xrightarrow{\unsure{e_{SA}|\varepsilon_{SA}|\psi(SA)}}(SA)$. Ceci dit, on a
\diagram{\begin{tikzcd}[ampersand replacement=\&]
 SA\arrow[r,"{S(\varphi(A))}"]\arrow[rr,bend right=17,"\text{identité}"']
 \& S(TS(A))=(ST)(SA)\arrow[r,"{\psi(SA)}"] \& SA
\end{tikzcd}}
en d'autres termes, $S(A)$ s'identifie par $S(\varphi(A))$ à un facteur direct de $(ST)(S(A))$, la projection naturelle étant $\deleted{e_{S(A)}}\psi(S(A))$. De \textquotedbl on a le diagramme commutatif
\diagram{\begin{tikzcd}[ampersand replacement=\&]
 T(\deleted{\unsure{A}}X) \& T(ST(\deleted{\unsure{A}}X))\deleted{\unsure{\leftarrow}}=\deleted{(TS)}(T(\deleted{\unsure{A}}X))\arrow[l,"{T(\psi(\deleted{\unsure{A}}X))}"'] \&
 T(X)\arrow[l,"{\varphi(T(\unsure{X|A}))}"']\arrow[ll,bend left=13,"\text{identité}"]
\end{tikzcd}}
qui identifie $T(X)$ à un facteur direct de $(TS)(T(X))$.

Pour que l'homomorphisme fonctoriel $\varphi:\mathrm{id}_C\longrightarrow TS$ soit injectif (i.e. les $\varphi(A):A\mathrel{-}TS(A)$ injectifs) il faut et il suffit que les homomorphismes
\[
 S(A,B):\mathrm{Hom}_C(A,B)\longrightarrow\mathrm{Hom}_{C'}(SA,SB)
\]
soient injectifs ; cela signifie aussi que le foncteur $S$ est \deleted{\unsure{fi}} ``injectif''. Cela résulte de la commutativité de
\diagram{\begin{tikzcd}[ampersand replacement=\&]
 \mathrm{Hom}_C(A,B)\arrow[r]\arrow[rr,bend right=12,"{\deleted{\unsure{u}}}"']
 \&\mathrm{Hom}_{C'}(SA,SB) \& \mathrm{Hom}_C(A,TSB)\arrow[l,"\sim"']
\end{tikzcd}}
% archive sheet number: 17
